\fsection{Ej 3 Pág 116}
\fsubsection
[\textcolor{waveBlue2}{TU SECTOR PROFESIONAL}]
{\textcolor{waveBlue2}{\boldit{TU SECTOR PROFESIONAL}}}
{
	\boldit{Tarea.}
	Nuevamente hemos trabajado sobre un ejemplo en el que el proceso de análisis de datos se aplica a una empresa de comercio
	electrónico, desde la comprensión del problema de negocio a partir de un objetivo claro hasta la presenta- ción de las
	recomendaciones para la empresa.\\

	En esta actividad, te invitamos a aplicar el proceso completo de aná- lisis de datos, o
	data analytics, a un negocio dentro de tu sector pro- fesional, en el que trabajes presencialmente con clientes y busques
	mejorar un aspecto del negocio. Anota en un documento todas las ideas asociadas a cada una de las fases y compártelo al final
	con tus compañeros y compañeras.\\

	Por ejemplo, en el sector sanitario, puedes imaginar una empresa (un laboratorio de análisis
	clínico, una parafarmacia, etc.) que busca mejorar la gestión de sus recursos y la atención al paciente. Comienza por
	comprender el problema de negocio definiendo claramente los objetivos; un posible ejemplo podría ser «Reducir los tiempos de
	espera en la sala de espera en un 20.
	Continúa con las cinco fases restantes aplicando cada una de ellas al ejemplo que hayas elegido.
}\\


\begin{solution}
	\textbf{Caso práctico: Optimización de afluencia y uso de equipamiento en centro de fuerza}

	\begin{itemize}
		\item \boldit{1. Comprensión del problema de negocio:}
		      \begin{itemize}
			      \item Objetivo claro: Reducir la saturación de las zonas de racks y peso libre en un 20\% durante las horas punta (18:00 - 20:00) para mejorar la retención de socios.
		      \end{itemize}

		\item \boldit{2. Recopilación de datos:}
		      \begin{itemize}
			      \item Registros de acceso mediante tornos (flujo de entrada/salida).
			      \item Datos de reservas de clases dirigidas (Powerlifting, Strongman) desde la App del centro.
			      \item Encuestas de satisfacción sobre disponibilidad de material.
		      \end{itemize}

		\item \boldit{3. Procesamiento y limpieza de datos:}
		      \begin{itemize}
			      \item Eliminación de registros duplicados en los accesos.
			      \item Normalización de tiempos de estancia (estimación de 90 min para salidas no marcadas).
			      \item Categorización por tipo de usuario (recreativo vs. competidor).
		      \end{itemize}

		\item \boldit{4. Análisis de datos:}
		      \begin{itemize}
			      \item Identificación de cuellos de botella: cruce de horarios de clases con picos de asistencia general.
			      \item Correlación entre el tiempo de espera percibido y la tasa de baja de socios (churn rate).
		      \end{itemize}
	\end{itemize}
\end{solution}
\newpage
\begin{solution}
	\begin{itemize}


		\item \boldit{5. Interpretación de resultados:}
		      \begin{itemize}
			      \item Se detecta que el 40\% de la saturación en racks proviene de clases grupales que coinciden con el horario de mayor acceso libre.
			      \item El segmento de competidores es el más afectado por la falta de disponibilidad de plataformas.
		      \end{itemize}

		\item \boldit{6. Presentación de recomendaciones:}
		      \begin{itemize}
			      \item \textbf{Gestión de horarios:} Desplazar las clases técnicas a horarios "valle" o zonas aisladas.
			      \item \textbf{Incentivos:} Implementar descuentos en cuota para usuarios que accedan en horario de mañana o mediodía.
			      \item \textbf{Monitorización:} Instalar un panel de ocupación en tiempo real en la recepción.
		      \end{itemize}
	\end{itemize}
\end{solution}

